Fundstellen werden in einer Literaturdatenbank gesammelt. Das Wort
"`Datenbank"' mag einige falsche Vorstellungen wecken. Für \LaTeX\ versteht man
darunter eine schlichte Textdatei im sog. Bib\TeX-Format.
Listing~\ref{lst:literatur:bibtexentry} zeigt einen beispielhaften Eintrag für
dieses Handbuch.

\lstinputlisting[language=BibTeX,style=block,float,caption={Ein BibTeX-Eintrag},label={lst:literatur:bibtexentry}]{code/literatur_bibtexentry.bib}

Ein Eintrag besteht aus drei Teilen.
%
\begin{compactenum}
\item dem Typ, hier \texttt{@manual}
\item einem Schlüssel, hier \texttt{urz\_vorlage\_manual}
\item mehreren Attributen, die die Fundstellen näher beschreiben
\end{compactenum}

Natürlich kann man die Datei manuell mit einem Texteditor pflegen. Es gibt
allerdings auch spezialisierte Software, die die Bearbeitung einer
Bib\TeX-Datenbank erleichtert. Als Beispiel wären zu nennen:
%
\begin{compactitem}
\item KBib\TeX~\cite{software_kbibtex}
\item Zotero~\cite{software_zotero}
\item JabRef~\cite{software_jabref}
\end{compactitem}
%
Mitunter speichern diese Programme die Literaturdatenbank in einem eigenen
Format und exportieren die Fundstellen nach der Bearbeitlung lediglich im
Bib\TeX-Format.

In den seltesten Fälle ist es notwendig, eine Bib\TeX-Datenbank händisch
anzulegen. Verlage und Bibliotheken halten in ihren Katalogen meist eine
entsprechende Exportfunktion vor, die es ermöglicht, vorgefertigte
Bib\TeX-Einträge herunterzuladen. I.d.R. sind jedoch anschließende Anpassungen
notwendig, falls Umlaute oder Sonderzeichen in den Autorennamen nicht korrekt
kodiert sind oder Datenfelder fehlen.
